% Part: first-order-logic
% Chapter: syntax-and-semantics
% Section: subformulas

\documentclass[../../../include/open-logic-section]{subfiles}

\begin{document}

\olfileid{fol}{syn}{sbf}

\olsection{\printtoken{P}{subformula}}

\begin{explain}
Kerep migunani kanggo ngrembug !!{formula}s kang nyusun sawijining
!!{formula}. Iki diarani \emph{!!{subformula}s} saka formula kasebut.
Saben !!{formula} dianggep !!a{subformula} saka awake dhewe;
subformula saka $!A$ kang dudu $!A$ dhewe diarani
\emph{!!{subformula} sejati}.
\end{explain}

\begin{defn}[!!^{subformula} Langsung]
Yen $!A$ iku !!a{formula}, \emph{!!{subformula}s langsung}
saka $!A$ ditetepake kanthi induktif kaya mangkene:
\begin{enumerate}
\item !!{formula}s atomik ora nduweni !!{subformula}s langsung.

\tagitem{prvNot}{\indcase{!A}{\lnot !B}{Siji-sijine
    !!{subformula} langsung saka $\indfrm$ yaiku~$!B$.}}{}

\item \indcase{!A}{(!B \ast !C)}{!!{subformula}s langsung saka
  $\indfrm$ yaiku $!B$ lan $!C$ ($\ast$ iku salah siji saka konektif
  rong panggonan).}

\tagitem{prvAll}{\indcase{!A}{\lforall[x][!B]}{Siji-sijine
    !!{subformula} langsung saka $\indfrm$ yaiku~$!B$.}}{}

\tagitem{prvEx}{\indcase{!A}{\lexists[x][!B]}{Siji-sijine
    !!{subformula} langsung saka $\indfrm$ yaiku~$!B$.}}{}
\end{enumerate}
\end{defn}

\begin{defn}[!!^{subformula} Sejati]
Yen $!A$ iku !!a{formula}, \emph{!!{subformula}s sejati}
saka $!A$ ditetepake kanthi rekursif kaya mangkene:
\begin{enumerate}
\item !!{formula}s atomik ora nduweni !!{subformula}s sejati.

\tagitem{prvNot}{\indcase{!A}{\lnot !B}{!!{subformula}s sejati saka
    $\indfrm$ yaiku~$!B$ bebarengan karo kabeh !!{subformula}s sejati
    saka~$!B$.}}{}

\item \indcase{!A}{(!B \ast !C)}{!!{subformula}s sejati saka
  $\indfrm$ yaiku $!B$, $!C$, bebarengan karo kabeh !!{subformula}s
  sejati saka $!B$ lan saka~$!C$.}

\tagitem{prvAll}{\indcase{!A}{\lforall[x][!B]}{!!{subformula}s
    sejati saka $\indfrm$ yaiku~$!B$ bebarengan karo kabeh
    !!{subformula}s sejati saka~$!B$.}}{}

\tagitem{prvEx}{\indcase{!A}{\lexists[x][!B]}{!!{subformula}s
    sejati saka $\indfrm$ yaiku~$!B$ bebarengan karo kabeh
    !!{subformula}s sejati saka~$!B$.}}{}
\end{enumerate}
\end{defn}

\begin{defn}[!!^{subformula}]
!!{subformula}s saka $!A$ yaiku $!A$ dhewe bebarengan karo kabeh
!!{subformula}s sejatine.
\end{defn}

\begin{explain}
Gatekna prabédan alus ing carane kita nemtokake !!{subformula}s
langsung lan !!{subformula}s sejati. Ing kasus kapisan, kita
langsung nemtokake !!{subformula}s langsung saka sawijining
formula~$!A$ kanggo saben wujud $!A$ kang mungkin. Iki definisi
eksplisit miturut kasus, lan kasus-kasuse niru definisi induktif
himpunan !!{formula}s. Ing kasus kapindho, kita uga niru cara
himpunan kabeh !!{formula}s ditetepake, nanging ing saben kasus kita
uga nyakup !!{subformula}s sejati saka !!{formula}s kang luwih cilik,
yaiku $!B$, $!C$, saliyane !!{formula}s kasebut dhewe. Iki ndadekake
definisi kasebut \emph{rekursif}. Umume, definisi sawijining fungsi
ing himpunan kang ditetepake kanthi induktif (ing kene, !!{formula}s)
iku rekursif yen kasus-kasus ing definisi fungsi mau nggunakake
fungsi kasebut dhewe. Nanging supaya definisine sah, kita kudu
mesthekake yen nilai fungsi kang digunakake mung kanggo argumen
kang teka luwih dhisik tinimbang argumen kang lagi ditetepake---ing
kene, nalika nemtokake "!!{subformula} sejati" kanggo $(!B \ast !C)$,
kita mung nggunakake !!{subformula}s sejati saka !!{formula}s
kang luwih dhisik, yaiku $!B$ lan $!C$.
\end{explain}

\begin{prop}
\ollabel{prop:subfrm-trans}
Upamana $!B$ iku subformula saka $!A$ lan $!C$ iku subformula saka $!B$.
Mula $!C$ iku subformula saka $!A$. Kanthi tembung liya, relasi
subformula iku transitif.
\end{prop}

\begin{prob}
Buktèkna \olref[fol][syn][sbf]{prop:subfrm-trans}.
\end{prob}

\begin{prop}
\ollabel{prop:count-subfrms}
Upamana $!A$ iku formula kanthi $n$ konektif lan kuantor.
Mula $!A$ nduweni paling akeh $2n+1$ subformula.
\end{prop}

\begin{prob}
Buktèkna \olref[fol][syn][sbf]{prop:count-subfrms}.
\end{prob}

\end{document}
